\documentclass[11pt,a4paper]{report}

\usepackage{lmodern}
\usepackage[T1]{fontenc}
\usepackage[utf8]{inputenc}
\usepackage[margin=2.6cm]{geometry}
\usepackage{amsmath}
\usepackage{booktabs}
\usepackage{graphicx}
\usepackage{microtype}
\usepackage{xcolor}
\usepackage{listings}
\usepackage[hidelinks]{hyperref}

\definecolor{codebg}{HTML}{F5F5F2}
\definecolor{rule}{HTML}{C3C2B7}
\lstset{
  basicstyle=\ttfamily\footnotesize,
  backgroundcolor=\color{codebg},
  frame=single,
  rulecolor=\color{rule},
  breaklines=true,
  columns=fullflexible,
  showstringspaces=false
}

% The same generated numbers the main report uses, so the two documents cannot
% disagree about a measurement.
% Generated by scripts/gen_report_assets.py. Do not edit by hand.
% Every value here comes from a file written by a run.
\newcommand{\AllGatesPass}{true}
\newcommand{\BenchCommit}{8e4c710216be}
\newcommand{\BenchFastMath}{on}
\newcommand{\BenchSeed}{20260802}
\newcommand{\BenchTimestamp}{2026-08-02T20:11:14Z}
\newcommand{\BestTileMs}{11.842}
\newcommand{\BestTileWidth}{256}
\newcommand{\CudaRuntime}{13.3}
\newcommand{\EnergyDriftMeasured}{$1.33 \times 10^{-5}$}
\newcommand{\FusedMsAtMaxN}{45.72}
\newcommand{\FusionCeilingMs}{0.0087}
\newcommand{\FusionCeilingPercent}{0.019}
\newcommand{\FusionGainPercentAtMaxN}{+0.1}
\newcommand{\FusionGainPercentAtMinN}{+2.1}
\newcommand{\FusionSpeedupAtMaxN}{1.001}
\newcommand{\GateCount}{10}
\newcommand{\GpuBandwidthEfficiency}{81.8}
\newcommand{\GpuBusBits}{192}
\newcommand{\GpuCompute}{12.0}
\newcommand{\GpuMeasuredBandwidth}{549.9}
\newcommand{\GpuName}{NVIDIA GeForce RTX 5070}
\newcommand{\GpuPeakBandwidth}{672.0}
\newcommand{\GpuSms}{48}
\newcommand{\LadderMaxN}{200,000}
\newcommand{\LadderMinN}{10,000}
\newcommand{\LsuReduction}{3.8}
\newcommand{\MeasuredFlopsPerInteraction}{18.0}
\newcommand{\MomentumDriftMeasured}{$1.23 \times 10^{-8}$}
\newcommand{\NaiveDramPct}{0.55}
\newcommand{\NaiveIssuePct}{50.1}
\newcommand{\NaiveLoneTexPct}{80.2}
\newcommand{\NaiveMsAtMaxN}{88.27}
\newcommand{\PeakTflops}{17.50}
\newcommand{\PlummerPValue}{0.4416}
\newcommand{\PolySixNormalisationError}{$6.48 \times 10^{-8}$}
\newcommand{\SectorReduction}{63.8}
\newcommand{\SpecifiedTileMs}{12.658}
\newcommand{\SphDensityErrorPercent}{4.20}
\newcommand{\SphFpsAtFiftyThousand}{2,963}
\newcommand{\SphFpsAtMaxN}{1,693}
\newcommand{\SphMaxN}{100,000}
\newcommand{\SphMsAtFiftyThousand}{0.337}
\newcommand{\SphMsAtMaxN}{0.591}
\newcommand{\TileWidthGainPercent}{6.4}
\newcommand{\TiledDramPct}{0.45}
\newcommand{\TiledIssuePct}{74.5}
\newcommand{\TiledLoneTexPct}{39.9}
\newcommand{\TiledMsAtMaxN}{45.78}
\newcommand{\TiledTflopsAtMaxN}{17.47}
\newcommand{\TilingSpeedupAtMaxN}{1.93}


% Each entry has the same shape: what I saw, what was actually wrong, what the
% options were, what I chose and why, and how I know it is fixed.
\newcommand{\entry}[1]{\section{#1}}
\newcommand{\field}[1]{\paragraph{#1}}

\title{\textbf{Particle Physics Simulation}\\[0.4em]
       \large Debug report: defects, surprises and corrections from the build}
\author{Olajide Badejo}
\date{August 2026}

\begin{document}

\maketitle

\begin{abstract}
\noindent
This is the companion to the main report and it ranks equally with it. The main
report says what the engine does; this one says what went wrong on the way and
what each failure taught, because in a project whose output is numbers the
history of the numbers that were wrong is worth as much as the ones that were
right.

Entries are grouped by theme. Environment covers the four toolchain defects
that had to be cleared before a single kernel compiled, one of which would have
silently produced a binary for the wrong GPU architecture. Performance covers
the two results that contradicted my expectations: fusion delivering
\FusionGainPercentAtMaxN{} percent where I expected ten to twenty, and shared
memory tiling turning out not to be a bandwidth optimisation at all. Physics
and tooling covers the defects in the correctness apparatus itself, including
two cases where the test rather than the code was wrong, and a data corruption
bug that put a random seed into a tile width column.

The most valuable entries here are the ones where nothing crashed. A build
targeting the wrong architecture, a benchmark reading a corrupted column and an
optimisation that quietly does nothing all produce plausible output, and none
of them announces itself.
\end{abstract}

\tableofcontents

\chapter{Environment}

Four defects stood between a clean checkout and a compiling kernel. None was
interesting in itself; all four are recorded because three of them produce
misleading rather than obvious failures.

\entry{nvcc 13.3 refuses GCC 16 as a host compiler}

\field{Symptom.} Compiling a trivial CUDA file with \texttt{-ccbin g++-16}:

\begin{lstlisting}
crt/host_config.h:137:2: error: #error -- unsupported GNU version! gcc
versions later than 15 are not supported!
\end{lstlisting}

\field{Root cause.} CUDA 13.3 supports GCC up to and including 15 and enforces
that with a preprocessor check, so this is a deliberate refusal rather than a
parse failure in some header. The machine has GCC 13, 14, 15 and 16 installed
and defaults to 15, so the problem only appears if the first design's floor of GCC
16.1 is taken literally, which it was.

\field{Options.} Pass \texttt{-allow-unsupported-compiler} and use GCC 16
anyway; drop to GCC 15.2, the newest supported host compiler; or install an
older toolkit.

\field{Fix and why.} GCC 15.2. What the floor was really buying is the C++23
host standard, and GCC 15.2 provides that in full. Suppressing the toolkit's
own compatibility check on a project whose entire output is numbers that are
supposed to be trustworthy is the wrong trade: a miscompiled force kernel that
still runs is the worst failure mode available here. Dropping the toolkit would
have given up first class \texttt{sm\_120} support on this card.

\field{Verification.} A kernel plus a CUB radix sort compiles, links and runs
correctly under \texttt{-ccbin g++-15 -arch=sm\_120 -std=c++23}.

\entry{CUDA 13 removed clockRate and memoryClockRate}

\field{Symptom.} The device probe would not compile:

\begin{lstlisting}
error: class "cudaDeviceProp" has no member "clockRate"
error: class "cudaDeviceProp" has no member "memoryClockRate"
\end{lstlisting}

\field{Root cause.} Both fields were deprecated during the CUDA 12 series and
removed in CUDA 13. The information is still available through
\texttt{cudaDeviceGetAttribute}.

\field{Fix and why.} Query both through the attribute interface. This matters
more than a compile fix normally would: the memory clock and bus width are what
the theoretical peak bandwidth is computed from, and the probe reports that figure beside
the measured bandwidth. Hard coding it from a data sheet would have meant quoting a number that had never
been measured on the machine that produced the results, which is exactly what
the no number without a run rule exists to prevent.

\field{Verification.} The probe reports \GpuBusBits{} bits at the measured
memory clock, giving \GpuPeakBandwidth{} GB/s, which agrees with the
manufacturer figure.

\entry{CMake could not identify the CUDA compiler at all}

\field{Symptom.} Configuration failed inside \texttt{project()}, before any
project code was reached:

\begin{lstlisting}
/usr/include/c++/15/type_traits(3303): error: __bool_constant is not a template
Error limit reached.
100 errors detected in the compilation of "CMakeCUDACompilerId.cu".
\end{lstlisting}

The confusing part was that nvcc had just compiled my own test kernel against
the same standard library without complaint.

\field{Root cause.} The difference was the standard flag. My test passed
\texttt{-std=c++20} explicitly. nvcc 13.3 defaults to C++17, and the libstdc++
15 headers it is paired with here use \texttt{if consteval} inside
\texttt{bits/c++config.h}, which is C++23 syntax. Under the default standard
the front end cannot parse its own standard library, so nvcc fails on an empty
translation unit. Reduced to two lines the failure is unambiguous:

\begin{lstlisting}
c++config.h(586): error: expected a "("
      if consteval { return true; } else { return false; }
\end{lstlisting}

CMake identifies the CUDA compiler by compiling a probe with
\texttt{CMAKE\_CUDA\_FLAGS} and nothing else. \texttt{CMAKE\_CUDA\_STANDARD} has
no effect there, because it applies to targets and no targets exist yet.

\field{Fix and why.} Seed \texttt{CMAKE\_CUDA\_FLAGS} with the standard before
\texttt{project()}, guarded so a value passed on the command line still wins.
Requiring every person who clones the repository to pass an extra flag would
have broken the promise that a clean clone plus two make commands reproduces
everything.

\field{Verification.} Configuration reports \texttt{The CUDA compiler
identification is NVIDIA 13.3.73 with host compiler GNU 15.2.0}.

\entry{The build silently targeted the wrong GPU architecture}

\field{Symptom.} None. Everything compiled and would have run. The only trace
was a status line reading \texttt{CUDA architectures: 75}, which is Turing, not
the Blackwell part in the machine.

\field{Root cause.} The default was being set with a guard placed after
\texttt{project()}:

\begin{lstlisting}
if(NOT DEFINED CMAKE_CUDA_ARCHITECTURES)
  set(CMAKE_CUDA_ARCHITECTURES 120)
endif()
\end{lstlisting}

Enabling the CUDA language inside \texttt{project()} assigns a default
architecture immediately, so by the time the guard runs the variable is already
defined and the branch never fires.

\field{Fix and why.} Move the guard above \texttt{project()} and make it a
cache variable, which is the only position where it can observe the undefined
state while still losing to a value passed on the command line.

The default is still \texttt{120} for development. \texttt{-DCPS\_FAT\_BINARY=ON} builds
\texttt{physics\_core} for six architectures, which is how the release and continuous
integration builds are made.

\field{Why this is the most important entry in the chapter.} Nothing broke.
The code would have run on the RTX 5070 through just in time compilation of the
embedded PTX, correctly and at the wrong speed. Every timing number in the
benchmark ladder would have been measured against a binary built for the wrong
architecture, and the only symptom would have been results somewhat worse than
expected, which is indistinguishable from the code simply being slow. A defect
that produces plausible numbers is far more dangerous than one that crashes.

% Two hyphens inside \texttt still form an en dash ligature; only a verbatim
% region suppresses it. The dash gate caught this one in the compiled PDF.
\field{Verification.} The status line now reads \texttt{CUDA architectures:
120}, and \verb|cuobjdump --list-elf| on a built object reports
\texttt{sm\_120} SASS rather than PTX alone. The latter check is the real one:
the status line reports intent, the object file reports fact.

\chapter{Performance}

Two expectations in the original design were contradicted by measurement, and the first
explanation of one of them was then contradicted as well. They are written down here
because an expectation that was wrong and got corrected is more useful than one that was
never recorded.

\entry{Kernel fusion, expected ten to twenty percent, measured about one}

\field{The expectation.} Fusing the force evaluation with the Verlet update takes one launch
instead of three and stops reading the acceleration array. Ten to twenty percent seemed
about right for a memory bound kernel.

\field{The measurement.} Round by round, at the default tile width, fused over tiled is
\FusedOverTiledRatioAtTenThousand{} at \LadderMinN{} particles and
\FusedOverTiledRatioAtHundredFiftyThousand{} at one hundred and fifty thousand, both beyond
their interquartile ranges and the session drift of \RepeatabilitySpread{} percent. At
\LadderMaxN{} particles it is \FusedOverTiledRatioAtMaxN{} with a range of
\FusedOverTiledRatioIqrAtMaxN{}, which does not resolve anything: the tiled trials are
bimodal, near \TiledLowModeMsAtMaxN{} and \TiledHighModeMsAtMaxN{} milliseconds. I do not
yet know why.

\field{Root cause of the wrong expectation.} Arithmetic intensity. The all pairs sum performs
$O(N^2)$ arithmetic on $O(N)$ data, so the ratio of operations to bytes grows without limit
with the problem size, and device memory is nearly idle: Nsight Compute puts it at
\NaiveDramPct{} percent of peak for the naive kernel and \TiledDramPct{} percent for the
tiled one.

\field{What I got wrong next.} \texttt{v1.0.0} turned this into a ceiling. It called the saving
a round trip of the acceleration array, $24N$ bytes written and read, priced that at the
measured bandwidth, and said no implementation could beat the result. Three things were wrong
with that. The fused kernel still wrote the array. The split path reads the array twice,
once in each integrator launch, so the description of the traffic was wrong too; counted from
the code, fusion removes \ModelBytesPerParticleSaved{}$N$ bytes, which is
\ModelSavingMsAtMaxN{} milliseconds or \ModelSavingPercentAtMaxN{} percent of a step at
\LadderMaxN{} particles. And the measured gains are several times larger than any such
figure, so it was never a ceiling. It prices one resource and says nothing about the rest of
the step. The gain itself, 0.1 percent in \texttt{v1.0.0}, was a difference of 0.12 percent
between separate blocks against a session spread of about 3 percent, so it could not be told
from noise.

\field{What I changed.} The protocol, so a difference of a percent can be read. The source of
the gain is still unmeasured. The fused kernel stays because it is faster where the difference
is resolved; \texttt{v1.0.0} also said it halves the memory footprint of a step, which is false,
because it adds a second position buffer.

\entry{Shared memory tiling is not a bandwidth optimisation on this problem}

\field{The expectation.} Tiling is conventionally introduced as a way of
reducing device memory traffic, and I introduced it that way. The prediction
that follows is that the naive kernel is limited by memory bandwidth and the
tiled kernel relieves it.

\field{The measurement.} Tiling delivers a factor of \TilingSpeedupAtMaxN{},
which is real and large. But device memory throughput is
\NaiveDramPct{} percent of peak before the change and \TiledDramPct{} percent
after. Device memory was never the constraint, so the conventional explanation
of why tiling helps is wrong here even though the tiling itself is right.

\field{What is actually happening.} The counters point at the first level cache
and the load store pipeline:

\begin{center}
\begin{tabular}{l rrr rr}
\toprule
& \multicolumn{3}{c}{throughput, percent of peak} & \multicolumn{2}{c}{per launch} \\
\cmidrule(lr){2-4}\cmidrule(lr){5-6}
variant & L1 and load store & L2 & device memory & global load sectors & warp issue \\
\midrule
naive & 87.98 & 1.62 & 0.34 & 5,000.1M & 54.99 \\
tiled & 45.41 & 1.61 & 0.35 & 78.3M & 82.68 \\
fused & 45.90 & 1.64 & 0.31 & 78.3M & 83.59 \\
\bottomrule
\end{tabular}

\end{center}

Global load sectors fall by a factor of \SectorReduction{} and load store
instructions by a factor of \LsuReduction{}. First level cache throughput falls
from \NaiveLoneTexPct{} percent of peak, close to saturation, to
\TiledLoneTexPct{} percent. The decisive number is warp issue efficiency, which
rises from \NaiveIssuePct{} percent to \TiledIssuePct{} percent. Tiling moves the reads into
shared memory and the load store pipe stops being the limit.

\field{Why this mattered practically.} Attacking the naive kernel's device memory traffic would
not have addressed the unit the counters name as limiting it. The profile pointed at the right
resource; reasoning from the familiar description would not have.

\field{Control.} The arithmetic instruction counts are equal across all three
variants to within a fraction of a percent, which is what licenses attributing
the difference to data movement rather than to the variants computing different
things. Those same counters also let the conventional
\ConventionalFlopsPerInteraction{} floating point operations per interaction be checked
against the \MeasuredFlopsPerInteraction{} actually retired, counting a fused multiply add as
two. Interactions per second is the measurement, and TFLOP/s depends on which count is used.

\entry{The first default tile width is not the fastest}

\field{Symptom.} A sweep of tile widths at \TileSweepParticles{} particles shows a width of
thirty two, the \texttt{v1.0.0} default, is well behind the widest of the four measured.

\begin{center}
\begin{tabular}{r rr rr}
\toprule
& \multicolumn{2}{c}{tiled} & \multicolumn{2}{c}{fused} \\
\cmidrule(lr){2-3}\cmidrule(lr){4-5}
tile width & ms per step & TFLOP/s & ms per step & TFLOP/s \\
\midrule
32 & 12.116 & 16.51 & 12.110 & 16.52 \\
64 & 11.777 & 16.98 & 11.769 & 16.99 \\
128 & 11.585 & 17.26 & 11.490 & 17.41 \\
256 & 11.487 & 17.41 & 11.373 & 17.59 \\
\bottomrule
\end{tabular}

\end{center}

\field{Root cause.} Two barriers are paid per tile, so a narrower tile pays them more often per
unit of arithmetic. That fits the direction of the result, though the barrier cost was not
measured on its own. A width of \BestTileWidth{} runs at \BestTileMs{} milliseconds against
\TileThirtyTwoMs{}, a difference of \TileWidthGainPercent{} percent, and a paired ratio of
\BestTileOverTileThirtyTwoRatio{} with an interquartile range of
\BestTileOverTileThirtyTwoRatioIqr{}.

\field{Fix and why.} \texttt{v1.0.0} kept thirty two and reported the sweep. The default is now
\BestTileWidth{}, and the ladder was measured there, because the difference is far outside the
session drift of \RepeatabilitySpread{} percent.

\field{A defect in the sweep itself.} The \texttt{v1.0.0} sweep timed the force kernel alone for
tiled, with particles that never moved, against whole steps for fused, so any comparison across
variants in its table was void. It also carried a claim here that thirty two was the slowest of
the four, when its own table had sixty four slower. Both variants now time a whole step.

\entry{Absolute timings depend on the thermal state of the card}

\field{Symptom.} In \texttt{v1.0.0}, a tile width sweep run immediately after the full benchmark
sweep produced timings about fifteen percent higher than the same sweep run from cold. Nothing
in the code had changed.

\field{Root cause.} Sustained clocks falling as the card heats is the likely reason, though I did
not log clocks. Either way it is large enough to swamp several of the effects being measured
here, including the entire fusion result.

\field{First fix, and why it was not enough.} \texttt{v1.0.0} compared within one session, kept the
fastest of three trials, and claimed relative differences were stable to a tenth of a percentage
point. They were not. Later measurement of the same binaries found them 4.7 percent apart in
separate blocks and 1.5 percent apart the other way when interleaved, because drift is confounded
with which variant ran when, and more trials do not help.

\field{Fix and why.} Measure in rounds: every configuration's first trial, then every
configuration's second, with the order rotating. Report the median of ten trials with its
dispersion. Take comparisons from the ratio inside each round, because two configurations in one
round saw the same machine. Measure one reference configuration at the start, middle and end of
the sweep, and treat the spread of those, \RepeatabilitySpread{} percent here, as the smallest
difference worth quoting.

\field{Verification.} Rerunning the sweep on a later commit reproduced the tile ratio to within a
tenth of a percent, and \texttt{scripts/check\_provenance.py} holds every recorded row to a commit
that resolves and a clean tree.

\chapter{Physics and tooling}

The entries here are about the correctness apparatus rather than the physics
kernels. Two of them are cases where the test was wrong, which is worth recording precisely because a wrong test is more dangerous than
a wrong implementation: it either blocks correct work or, far worse, passes
when it should not.

\entry{The hash collision that would have counted neighbours twice}

\field{The pathology.} Spatial hashing is many to one by construction. Two
distinct grid cells can map to the same bucket. The neighbour search visits the
twenty seven cells around a particle and scans each one's bucket, so if two of
those twenty seven cells collide in the table, the particles in that bucket are
scanned twice and counted twice.

The reason this is nasty is that the usual defence does not work. Filtering
candidates by distance removes particles from unrelated cells, because they are
genuinely far away. It does nothing about double counting, because those
particles really are within the smoothing radius; they are simply being added
to the sum more than once. The density would come out high by an amount that
depends on the hash function rather than on the fluid, and the fluid would
still look like a fluid.

\field{Fix and why.} Each particle stores an exact cell identity packed into 64
bits as three 21 bit fields, and the bucket scan compares it before accepting a
candidate. This rejects both classes at once, at the cost of one 64 bit compare
per candidate. The alternative of deduplicating the twenty seven hashed buckets
before scanning would cost hundreds of comparisons per particle rather than
one per candidate.

\field{Verification.} This is the one defect in the project that was designed
against rather than discovered, so the verification carries the whole weight.
The density field produced by the hashed search is compared against a brute
force $O(N^2)$ host sum over identical positions, and the worst relative
difference across four thousand particles is \NeighbourSearchError{}, about what single
precision rounding gives. A single dropped or double counted neighbour
would move a density by several percent. A separate health check counts
particles sharing a bucket with a foreign cell and confirms the collision rate
is real but small, so the test is exercising the defence rather than a table so
sparse that no collisions occur.

\entry{A diagnostic that would have perturbed the run it measured}

\field{The problem.} The fused variant carries the leapfrog half step velocity,
which is not the velocity belonging to the current positions. The energy
diagnostic needs the synchronised one. The obvious implementation is to shift
the velocity array forward by half a step, take the measurement, and shift it
back.

\field{Why that is wrong.} In single precision, adding a quantity and then
subtracting it does not return the original bits. A run that happened to take a
diagnostic would therefore follow a very slightly different trajectory from one
that did not, which breaks the bit identity guarantee the resume path claims
and makes the energy curve partly a measurement of its own observer.

\field{Fix.} The half step correction is applied inside the reduction kernel,
which reads the velocity and the stored acceleration and forms $v + a\,\Delta
t/2$ on the fly without writing anything back.

\field{Verification.} A test runs eighty steps taking a full diagnostic after
every step, and asserts the final state is bit for bit identical to the same
run with no diagnostics at all.

\entry{The density gate that passed by measuring the wrong particles}

\field{Symptom.} None, which is the point. \texttt{v1.0.0} published an interior density error of
4.20 percent against a 5 percent limit, with an explanation about smoothing on a jittered lattice.
It was found by review for this release rather than by any failure.

\field{Root cause.} Two defects, one hiding the other. The gate chose interior particles by position
in lattice order and read densities in cell sorted order, so positions and densities belonged to
different particles, and the figures were close to the statistics of the whole block. With the index
fixed, the \texttt{v1.0.0} mass gave 7.61 percent, which fails that limit. And that mass was wrong: it
assumed a spacing the lattice does not have, and ignored the kernel's own lattice sum error.

\field{Consequence.} The equation of state turns excess density into pressure, so the
\texttt{v1.0.0} dam break interior started near 76 kPa against a hydrostatic 3.9 kPa. Maximum speed
after 100 steps was 3.695 m/s. With the calibrated mass it is 0.625 m/s.

\field{Fix.} The index is fixed, the mass comes from a kernel sum over the lattice laid out, and new
limits were committed before the code measuring them existed: the mean error is
\SphDensityErrorPercent{} percent, the RMS deviation \SphDensityRmsPercent{} percent and the maximum
deviation \SphDensityMaxPercent{} percent.

\field{Why it is in this report.} The gate read well, passed, and came with a plausible story. A
check that was never seen to fail for the right reason should not be trusted just because it passes.

\entry{A test that failed for a reason that was entirely my fault}

\field{Symptom.} The SPH kernel unit test
\texttt{KernelsAreContinuousAtTheSupportBoundary} failed on first run, while
every other kernel test passed.

\field{Root cause.} The test asserted that each kernel evaluated just inside its
support boundary was within $10^{-3}$ of zero, in absolute terms. These kernels
carry a $1/h^6$ normalisation: at a smoothing radius of $0.1$ the spiky
gradient peaks around $1.4 \times 10^5$ and the viscosity Laplacian around
$1.4 \times 10^6$. An absolute tolerance of $10^{-3}$ against quantities of that
magnitude is not a statement about continuity, it is a statement about nothing.
The implementation was correct throughout.

\field{Fix.} The test now compares each kernel at the boundary against its own
peak value, which is what continuity actually means for a quantity with a
scale, and separately asserts the kernels are exactly zero outside the support
so the transition is genuinely continuous rather than merely small.

\field{Lesson.} A tolerance without a scale is not a tolerance. Every numerical
gate in this project is now expressed relative to something: the magnitude of
the quantity, the machine epsilon, or a bound derived from the arithmetic.

\entry{A benchmark that recorded a random seed as a tile width}

\field{Symptom.} The tile width figure had four bars, one of which sat at a
tile width of \texttt{20260802}. That number is the random seed.

\field{Root cause.} The tile sweep opened the output file twice. It created a
truncating output stream to write the header, and then called the row appending
helper, which opens the same path independently in append mode. The header
stream stayed open and buffered for the duration of the sweep and flushed at
offset zero when it was finally destroyed, overwriting the beginning of the
file with the header text. What survived was a file whose first rows were
partly overwritten, so the column boundaries shifted and fields were read from
the wrong columns.

\field{Fix.} The header is written through the same append path as the rows,
with the file removed first, so only one handle to the file exists at any time.

\field{Why it is in this report.} Nothing failed. The benchmark ran, the CSV
parsed, the plot rendered, and the number was wrong. It was caught by looking
at the figure and asking what a tile width of twenty million meant. The general
form of this defect is that two independent handles to one file will corrupt
it, and buffering makes the corruption appear far from its cause.

\entry{LaTeX macro names cannot contain digits}

\field{Symptom.} The report failed to compile with
\texttt{Missing \textbackslash begin\{document\}} pointing at a line in the
generated macro file that looked entirely reasonable.

\field{Root cause.} A LaTeX control sequence is made of letters only. The
offending definition was

\begin{lstlisting}
\newcommand{\Poly6NormalisationError}{...}
\end{lstlisting}

\noindent
in which the name ends at the digit, so this defined a macro called
\texttt{Poly} followed by the literal text \texttt{6NormalisationError}. Literal
text in the preamble is what produces that error message, and it is reported at
the point of the text rather than at the mistake.

\field{Fix.} The three offending macros were renamed to spell their digits, and
the generator now rejects any macro name that is not purely alphabetic, so this
class of mistake fails at generation time with a message that names the
problem rather than at compile time with one that does not.

\entry{BibTeX put en dashes back after they had been removed}

\field{Symptom.} The dash gate failed on the compiled PDF, reporting seven en
dashes in the bibliography, after the source bibliography had been cleaned of
double hyphens and the gate had reported the sources clean.

\field{Root cause.} The bibliography style is what dashifies page ranges. The
\texttt{plain} style passes page fields through a function that converts any
run of hyphens into a TeX en dash ligature when it writes the intermediate
bibliography file. Cleaning the source had no effect because the source was
never the problem.

\field{Fix.} Page ranges are written with the word \texttt{to}, which the style
passes through untouched. This also makes the bibliography consistent with the
prose, which is required to write ranges the same way.

\field{Why this justifies the whole design of the gate.} The rule forbids a
character in the finished document, so the gate has to check the finished
document. A source only check would have reported clean while every reference
in the printed bibliography broke the rule. Extracting the text of the compiled
PDF and scanning that is what caught it, and it is the reason the gate runs
last in the report build rather than first.

\chapter{What the failures had in common}

\section{The dangerous defects were the silent ones}

Of the eleven entries in this report, three crashed, three failed to compile,
and five produced plausible output while being wrong. The five silent ones took
almost all the time.

A build targeting Turing instead of Blackwell compiles, links and runs. A
benchmark whose columns have shifted parses cleanly and plots. An optimisation
that relieves a resource nothing was waiting on delivers a small positive
number that is easy to accept. A neighbour search that counts some particles
twice produces a fluid that still behaves like a fluid. A diagnostic that
perturbs the state it measures produces a smooth curve.

None of these announces itself, and none is found by testing that the program
runs. Each was caught by a different mechanism, and the pattern is that each
mechanism compares the thing against something independent: the object file
against the intended architecture, the figure against what a tile width could
plausibly be, the optimisation against a bound computed from measured
bandwidth, the neighbour search against a brute force sum, the diagnostic run
against an undiagnosed one.

\section{A tolerance without a scale is not a tolerance}

Two defects came from this directly. The kernel continuity test compared a
quantity of order $10^6$ against an absolute tolerance of $10^{-3}$ and failed a
correct implementation. Had the comparison gone the other way, against a
quantity of order $10^{-6}$, it would have passed anything at all.

Every gate in the project is now expressed against a scale that means
something. The momentum gate is bounded by $\epsilon_{\mathrm{f32}} \sqrt{N}
n_{\mathrm{steps}}$, which is what single precision arithmetic can explain, and
it passes with six orders of magnitude to spare, so the number reported is
informative rather than merely a pass. The force gate is relative to the
magnitude of the acceleration being checked. The density gate is relative to
rest density.

\section{Size the optimisation before writing it}

The fusion result is the clearest lesson in the project. The saving was
computable in advance from two numbers already measured at Phase 0: the bytes
removed and the achievable bandwidth. That calculation gives
\FusionCeilingPercent{} percent as an absolute ceiling at \LadderMaxN{}
particles and takes two minutes.

The same calculation would have reframed the tiling work correctly as well.
Tiling is worth a factor of \TilingSpeedupAtMaxN{}, which no bandwidth argument
can explain, because bandwidth was never the constraint. Knowing that in
advance would have pointed at the load store pipeline from the beginning
instead of at the memory bus.

\section{Verification has to be independent of the thing verified}

The strongest checks in this project are the ones where the reference was
produced by a genuinely different route. The double precision host force
implementation shares no code with the kernels and is written as literally as
the force law can be written. The brute force density sum shares no code with
the neighbour search. The chi squared test is checked by feeding it a
deliberately wrong model that it must reject, so a test that could never fail
would itself fail.

Where sharing is unavoidable it is deliberate and in the safe direction: the
smoothing kernels and the force law are compiled from one header into both the
device kernels and the host tests, so the two cannot drift apart. That is
sharing an implementation to guarantee agreement, not sharing an assumption to
guarantee a pass.

\section{The environment is part of the experiment}

Four of the eleven entries were toolchain problems, and one of them would have
invalidated every performance number in the main report while leaving no trace
except numbers that were slightly disappointing. Phase 0 of this build existed
to probe the environment and record what it found, and that turned out to be
time well spent rather than ceremony. The measured bandwidth figure recorded
there is the denominator of the fusion ceiling calculation that explains the
central performance result of the project.


\end{document}
